\documentclass[a4paper,11pt]{article}
\usepackage[utf8]{inputenc}
\usepackage[T1]{fontenc}
\usepackage[english]{babel}
\usepackage[left=2.5cm,right=2.5cm,top=2cm,bottom=2cm]{geometry}
\usepackage{hyperref}
\usepackage{xcolor}
\usepackage{microtype}
\usepackage{lmodern}

\definecolor{primary}{RGB}{0, 51, 102}
\hypersetup{colorlinks=true, urlcolor=primary}

\begin{document}
\noindent
\textbf{Lo\"ic Aron MBASSI EWOLO} \\
ENSEA -- Double Dipl\^ome Ing\'enieur \\
\href{mailto:loicaron.mbassiewolo@ensea.fr}{loicaron.mbassiewolo@ensea.fr} \quad (+33) 7 59 52 33 98 \\
\href{https://github.com/Nameless0l}{github.com/Nameless0l} \quad \href{https://mbassiloic.tech}{mbassiloic.tech}

\vspace{1em}
\noindent
Cergy, le \today

\vspace{1em}
\noindent
\textbf{Objet : Candidature -- Data Scientist IA} \\
\textbf{R\'ef\'erence : Offre d'alternance 12 mois -- BPCE}

\vspace{1.5em}
\noindent Madame, Monsieur,

\vspace{0.8em}
\noindent En tant que groupe mutualiste de premier plan porteur d'innovation financière, BPCE/Banque Populaire poursuivez une transformation numérique ambitieuse appuyée sur l'intelligence artificielle et la data science. Face aux enjeux de scoring de risque, détection fraude et recommandations produits, l'IA générative et les pipelines NLP innovants représentent des leviers stratégiques. C'est pourquoi je candidats avec enthousiasme au poste de Data Scientist IA en alternance dès septembre 2026, convaincu que mon expérience INRIA en NLP production et mon engagement pour l'IA responsable correspondent exactement à votre vision.

\vspace{0.8em}
\noindent Depuis avril 2026, je déploie un pipeline NLP complet au sein de l'équipe TRIBE (INRIA Saclay) : fine-tuning BERT sur corpus de politiques de confidentialité, classification multi-labels pour scoring de conformité RGPD, et agrégation d'un score interprétable (0-100) de confiance. Cette expérience directe m'a particulièrement formé aux enjeux d'une IA en production bancaire : latence < 100ms par requête, intégrabilité dans des systèmes de décision existants, et explicabilité des résultats pour la conformité réglementaire. J'ai approfondi l'analyse de risques (détection SDKs malveillants exfiltrant données), ce qui transpose directement aux problématiques BPCE de détection fraude et évaluation de risques crédit.

\vspace{0.8em}
\noindent En complément, j'ai diriger et architecter un système d'agents IA multi-agents orches-trés par Google ADK (Planning, Météo, Marché, Recommandations), combinant LLM Gemini avec une RAG enrichie d'APIs externes. Cette expérience de coordination d'agents autonomes, de résolution de conflits multi-agents et de recommandations multi-critères ilustre ma capacité à penser systémiquement des architectures IA complexes. Ma recherche chez MILA (Montréal) sur le phénomène de grokking m'a instillé la rigueur scientifique indispensable pour valider la robustesse des modèles, avec tests statistiques rigoureux et analyse de dynami-ques d'apprentissage. J'affirme que cette combinaison (NLP production + agents IA générative + rigueur research) positionne parfaitement pour contribuer à l'innovation IA de BPCE.

\vspace{0.8em}
\noindent Je suis disponible dès septembre 2026 pour une alternance 12 mois, ce qui m'permettra de déployer des initiatives data science concrets au service de vos utilisateurs et de vos équipes métiers. Mes expériences en scoring de risques, détection d'anomalies et recommandations intelligentes seront immédiatement transférables aux challenges banc-aires. Je serais ravi de discuter comment mon expertise INRIA en NLP, mes compétences en IA générative et agents, et ma vision d'une IA responsable et explica-ble peuvent servir l'ambition numérique de BPCE. Merci pour cette opportunité.

\vspace{1.5em}
\noindent Veuillez agr\'eer, Madame, Monsieur, l'expression de mes salutations distingu\'ees.

\vspace{2em}
\noindent \textbf{Lo\"ic Aron MBASSI EWOLO}
\end{document}
